\chapter{The HTTP semantic model}
\chapterintro{HTTP is a family of stateless request/response protocols with shared semantics. First learn what a message means; the next chapter studies exactly how HTTP/1.1 writes it as bytes.}

\begin{objectives}
\begin{itemize}
  \item Distinguish URI, origin, resource, representation, content, and field metadata.
  \item Explain method safety, idempotency, and the major method semantics.
  \item Interpret status-code classes without treating every 2xx or 4xx alike.
  \item Explain stateless interaction and where application state still lives.
\end{itemize}
\end{objectives}

\section{Messages carry intentions and results}

An HTTP client sends a request that identifies a target and expresses an
intention through a method, fields, and optional content. A server interprets
that request with respect to a resource and returns one or more responses. A
final response has a status code, fields, and optional content.

HTTP separates a uniform interface from implementation details. A client does
not need to know whether a task is stored in SQLite, computed from another
service, or generated temporarily. The server does not need to know why the
client wants it. This separation allows both sides to evolve behind the protocol
boundary.

\section{URI, origin, resource, representation}

A URI identifies a resource. For HTTP and HTTPS, an \term{origin} is the tuple
of scheme, host, and port after applying defaults. Thus
\url{http://example.test/tasks} and
\url{https://example.test/tasks} are different origins, even though their host
and path text match.

The resource is conceptual. Its current representation is a sequence of content
bytes plus metadata describing how to interpret them. A JSON media type tells a
recipient to process bytes as JSON; a character encoding connects text bytes to
characters. Representation metadata can also describe language, compression,
freshness, and validators.

\section{Methods are semantics, not route names}

The method token is case-sensitive and states the request's primary semantics.
Standard method names are uppercase.

\begin{center}
\begin{tabularx}{\textwidth}{@{}>{\bfseries}p{20mm} p{23mm} p{25mm} X@{}}
\toprule
Method & Safe & Idempotent & Intended meaning \\
\midrule
GET & yes & yes & Transfer a current representation \\
HEAD & yes & yes & Same as GET without response content \\
POST & no & no & Perform resource-specific processing \\
PUT & no & yes & Replace the target's current representation \\
PATCH & no & not guaranteed & Apply a described partial modification \\
DELETE & no & yes & Remove the association with the target resource \\
OPTIONS & yes & yes & Describe communication options \\
\bottomrule
\end{tabularx}
\end{center}

\term{Safe} means the requested semantics are read-only; logging and accounting
side effects can still occur. \term{Idempotent} means repeating the same request
has the same intended effect as performing it once. It does not require byte-for-
byte identical responses. A repeated DELETE can first return success and later
return not found while the resource remains deleted.

These properties guide retries. A client can often retry an idempotent operation
after losing a connection without fearing duplicate intended effects. Retrying a
POST that creates a charge or task may duplicate work unless the API adds an
idempotency mechanism.

\begin{pitfall}
Using GET for a state-changing action because it is convenient violates method
semantics. Crawlers, caches, link previewers, and retrying clients are entitled
to treat GET as safe.
\end{pitfall}

\section{Status codes describe the result}

The first digit groups status codes:

\begin{itemize}
  \item \textbf{1xx Informational}: an interim response before a final response;
  \item \textbf{2xx Successful}: the request was received, understood, and accepted;
  \item \textbf{3xx Redirection}: another step or a cached representation is involved;
  \item \textbf{4xx Client Error}: the request cannot be fulfilled as sent;
  \item \textbf{5xx Server Error}: the server failed to fulfill an apparently valid request.
\end{itemize}

Choose the specific code that communicates the outcome. \status{200 OK} commonly
carries a representation. \status{201 Created} identifies a newly created
resource, usually with a \code{Location} field. \status{204 No Content} reports
success without content. \status{400 Bad Request} means malformed or invalid
request syntax at a broad level; \status{404 Not Found} hides or reports absence;
\status{405 Method Not Allowed} means the resource exists for other methods and
should include \code{Allow}. \status{415 Unsupported Media Type} rejects the
content format. \status{422 Unprocessable Content} can report semantically invalid
fields. \status{500 Internal Server Error} is not a license to reveal a traceback.

\section{Fields modify or describe messages}

HTTP fields are named metadata. Names are case-insensitive; values follow
field-specific syntax. Some fields can appear more than once or contain lists,
so blindly storing every name in a one-value dictionary loses information.

Examples include:

\begin{itemize}
  \item \code{Host}, which identifies the target authority in HTTP/1.1;
  \item \code{Content-Type}, which describes content media type;
  \item \code{Content-Length}, which participates in HTTP/1.1 framing;
  \item \code{Accept}, which expresses response representation preferences;
  \item \code{Authorization}, which carries credentials under a chosen scheme;
  \item \code{Cache-Control}, which gives cache directives;
  \item \code{Location}, which identifies a created or redirect target.
\end{itemize}

Fields are not all interchangeable strings. Their individual semantics determine
whether whitespace, combination, repetition, and ordering matter.

\section{Stateless does not mean no state}

HTTP is stateless in the sense that each request can be understood without an
implicit transport conversation state defining its meaning. Application state
still exists in databases, caches, and clients. A client can send a cookie or
authorization credential on each relevant request, making context explicit in
the request.

Persistent TCP connections are a performance mechanism, not an application
session. Two requests on one connection should not secretly mean different
things merely because of their connection order, unless a protocol extension
explicitly defines such state.

\section{Content versus representation}

Modern HTTP specifications use \term{content} for the bytes after transfer
codings have been removed. In HTTP/1.1, chunk boundaries are message framing;
they are not part of the content delivered to an application. Likewise, TLS
records protect transport bytes but are not part of the HTTP representation.

This distinction keeps semantics consistent across versions. HTTP/2 uses binary
frames and HTTP/3 uses QUIC streams, but a GET, status 404, media type, and cache
directive retain shared meaning.

\summarytitle

HTTP gives resources a uniform request/response interface. Methods express
intentions and properties such as safety and idempotency. Status codes communicate
results. Fields carry typed metadata. Representations transfer resource state,
while HTTP's stateless model keeps each request independently meaningful.

\reviewtitle
\begin{enumerate}
  \item Why are a resource and its JSON representation different things?
  \item How can an idempotent request return different responses when repeated?
  \item What distinguishes 404 from 405?
  \item Why is a persistent TCP connection not an application session?
\end{enumerate}

\begin{exercise}[Method audit]
Choose methods for creating, reading, replacing, partially updating, and deleting
a task. For each, state whether an automatic retry is generally safe and which
application-specific caveat remains.
\end{exercise}

\begin{exercise}[Choose the result]
Choose and justify status codes for: successful creation, successful deletion
with no body, malformed JSON, an unknown task, a known path with the wrong
method, an unsupported media type, invalid task fields, and an unexpected
database failure.
\end{exercise}
